\section{Conclusions and Future Improvements}
\label{sec:conclusions}

\subsection{What the measurements say}

\textbf{Strengths.} The design is correct under every failure we could cause. All 165 recovery runs recover with the
same answers whatever the mode and the replica count, and no book ever appeared twice. Coordination costs little: a
message takes 0.1--0.2\,ms, and the broker uses at most 17\,\% of one core. Indexer and \code{query} replicas raise
throughput on one host, $\times$2.10 with three indexers and +29\,\% with a second \code{query}. A book becomes
searchable about 40\,ms after its batch is published. The contracts held: the API checks pass in every mode and
with one or three \code{query} replicas, and every figure traces to a run.

\textbf{Weaknesses.} Indexing is the bottleneck: one indexer publishes 28.5 books/s against about 500 downloaded, and
its time goes into MongoDB writes. Crawler replicas add nothing beyond 500--670 books/s, because they share one queue,
one database and one volume. \code{control} is a single process that every mode needs, since it alone releases leases
and drives rebuilds. When the broker loses messages, \code{control} brings them back after
\code{TARANTINO_REQUEUE_SECONDS}, slower than the redelivery of a persistent broker. One open question is why that
requeue took 82--97\,s after a Hazelcast restart with 2 and 3 replicas, against 12\,s with one.

\subsection{From one host to the Stage~3 cluster}
\label{sec:stage3}

The conditional writes already work across machines. Claims, leases and expirations compare with MongoDB's clock,
never a host's, and no process holds a lock that another needs. What Stage~3 has to replace is everything named after
the host (Table~\ref{tab:stage3}): Compose service names and nginx's resolver \code{127.0.0.11}, the local
\code{datalake} volume, the single Hazelcast member and MongoDB server, and, in mode \mode{orchestrator}, the JVM's DNS
cache through which \code{control} finds the \code{query} replicas.

\begin{table}[ht]
\centering
\small
\caption{Each measured limit, the single-host assumption behind it, and the change Stage~3 makes.}
\label{tab:stage3}
\begin{tabularx}{\textwidth}{@{}L{4.0cm}L{4.1cm}Y@{}}
\toprule
Measured limit & Single-host cause & Stage~3 change \\
\midrule
Indexing: 28.5 books/s per indexer, $\times$2.10 with 3, with MongoDB's CPU growing faster than the indexers' & Every indexer writes postings to one MongoDB server & Postings sharded by \code{term} across a sharded MongoDB, so indexers on several nodes write to different shards; the batch, its lease and its publication stay single-document writes \\
Ingestion: 500--670 books/s whatever the replicas, with the CPU far from full & One queue on one Hazelcast member, one MongoDB, one datalake volume & A Hazelcast cluster with the queue partitioned by \code{book_id}, and a datalake on shared storage, or one per node with the book's node in \code{book_states} \\
Search: one \code{query} replica saturates at about 4{,}400 requests/s, and a third adds nothing to two & Replicas share nginx, the one MongoDB and the host's CPUs with the load generator & \code{query} replicas on several nodes behind a load balancer, reading a MongoDB replica set; the snapshot already lets each replica answer on its own \\
Recovery: lost messages wait for \code{control}'s requeue & \code{control} is one process, acceptable on one host & \code{control} on several nodes with one leader, elected through a lease document in MongoDB; lost work acknowledged after the job (below) \\
Mode \mode{orchestrator}: 41\,\% less throughput with 3 replicas & One thread dispatches every book, one \code{POST} at a time & Not carried to Stage~3 \\
\bottomrule
\end{tabularx}
\end{table}

\textbf{Communication and storage across machines.} Services keep talking only through the broker and MongoDB, so
moving them to other nodes changes addresses, not code. Hazelcast becomes a cluster of three members, where each queue
item has a synchronous backup on another member, so one member can fail without losing queued work. MongoDB becomes a
replica set, with majority writes for every conditional write, so a failover never loses a claim. The datalake moves to
storage that every node mounts.

\textbf{Distribution of tasks and requests.} Crawlers and indexers keep competing for queues, which needs no
coordination between them. Partitioning \code{books-to-download} and \code{books-to-index} by \code{book_id} removes
the single queue owner that limits ingestion today. Search requests are spread by a load balancer, as nginx does now.

\textbf{Consistency and unavailable nodes.} Publication stays one conditional write, and \code{query} keeps reading a
snapshot of published batches, so a reader never sees a half-built index on any node. A node that disappears leaves
leases that expire, and \code{control} releases them exactly as it does after a crash today: rows 01 to 10 of the
failure matrix already exercise that path. A network partition cannot create two owners of a book, because ownership
is a MongoDB document, not a broker message.

\subsection{The winning mode in Stage~3}

\mode{hazelcast} wins under the decision rule fixed before the campaign (Section~\ref{sec:winner}), so Stage~3 keeps it
and does not migrate to another broker. The runner-up, \mode{activemq-persistent}, showed the one property worth
taking: lost work comes back without \code{control}, in 3.8--6.1\,s instead of about 12\,s, because the message is
acknowledged only after the job. Stage~3 adds that property to the Hazelcast design and measures it with the same
rows and the same \code{recovery-delay} experiment:
\begin{enumerate}
\item Take each book inside a Hazelcast transaction that commits after the claim and the job, so a consumer that dies
  returns its message to the queue. Hazelcast's transactional queue offers that rollback~\cite{hazelcast}, and the claim in MongoDB
  still drops any duplicate it produces.
\item Explain and fix the 82--97\,s requeue after a broker restart with several replicas (row 12).
\item Elect the \code{control} leader through a lease in MongoDB, so reconciliation survives the loss of a node.
\item Repeat \code{coordination-variants} for \mode{hazelcast} and \mode{activemq-persistent} on the cluster. If the
  transactional queue does not close the gap in rows 04, 05 and 12, persistent ActiveMQ replaces Hazelcast as the
  broker. The rest of the code stays the same, because both already sit behind the same ports.
\end{enumerate}

Smaller improvements follow from this report: batches of 250 books once the indexer has a larger heap, which gave
17\,\% more throughput, and ``Open a rebuild'' computing its target inside its own write, which closes the rebuild race
of Section~\ref{sec:indexer}.

\subsection{Lessons}

\begin{itemize}
\item \textbf{Measure the alternatives.} Arguing against orchestration and polling was easy; measuring them showed
  where each one actually loses, by how much, and what the runner-up does better.
\item \textbf{Keep the truth in one place.} Treating every message as a hint and every transition as one conditional
  write made 14 different failures recoverable by the same reconciliation.
\item \textbf{Keep the evidence with the results.} Five interleaved repetitions, hypotheses written before measuring,
  and a versioned run behind every figure made every number in this report checkable against its raw data, and we
  checked them all before submitting it.
\end{itemize}
